\scanpage{194}
(fonctions \deleted{symétriques}\inserted{invariantes} par l'action du groupe
$\mathfrak S_{\{i,j,k\}}\subset\mathfrak S_4$).
\unsure{Ce fait|Ça fait} qu'en fait, il suffit comme générateurs les quatre quantités
(en plus des quatorze $b_\cdot,c_\cdot,t_{\cdot\cdot}$ du début)
\[
t_{234},\quad t_{1,3,4},\quad t_{1,2,4},\quad t_{1,2,3}.
\]
Si $A$ est l'anneau engendré par les $b_\cdot,c_\cdot,t_{\cdot\cdot}$,
je présume que $B$ est fini libre de rang $2^4$ sur $A$, \unsure{aux|avec}
$8$ gén.\inserted{distincts $t_{ijk}$} sur $A$ soumis aux seules relations $(*)$ ci-dessus.

Mais il faudrait trouver la relation entre les $b_i,c_j,t_{k,l}$ --- peut-être via les
$t_{ijk}$. La clef pour la trouver \deleted{est}\inserted{sera} le calcul de l'expression
\[
\operatorname{Tr}x_1x_2x_3x_4
\]
\margin{193}
